

\chapter{Decisiones de Diseño}
 
\section{Paradigmas y estrategia general}
 
Para el diseño de los algoritmos paralelos se adoptaron dos paradigmas de forma
complementaria:
 
\begin{itemize}
  \item \textbf{Paradigma de datos} para la distribución de los datos de entrada: dentro de
        cada tarea, los datos (filas o bloques de la imagen) se reparten entre los workers
        disponibles.
\end{itemize}
 
Dentro de la sección de borroneado, las dos subtareas (blur y highlight+umbralado) sí
presentan dependencia secuencial entre sí, por lo que se implementa un \textbf{pipeline de
datos}: primero se ejecuta la función de blur sobre toda la imagen y, sobre el resultado,
se aplica la función de highlight junto con la operación de umbralado.
 
\section{Algoritmo secuencial}
 
En todas las versiones del algoritmo se aplicó la técnica de \textbf{padding}: se añade un
borde auxiliar de píxeles con valor neutro (negro) alrededor de la matriz de entrada. Esto
elimina la necesidad de verificar en cada iteración si las celdas vecinas son válidas al
calcular el blur y el highlight, simplificando el código y mejorando la localidad de memoria.
 
\section{Algoritmo implementado en CUDA}

